
The Bridge Rapid Assessment Center for Extreme Events (\BRACE{}) is a platform developed in collaboration with the California Department of Transportation to provide real-time post-earthquake assessments for instrumented bridges.
Its computational backend, the Infrastructure Resilience Engine (IRiE), requires finite element models that can be assembled from inventory data, executed rapidly, and embedded in inference workflows.
The bridges in the current deployment are composed predominantly of beam and column members whose inelastic behavior governs damage progression.
The beam formulations developed here serve as the primary simulation engine, and the modular structure of the framework---permitting, for instance, a transition from a coarse linear model to a refined inelastic model by changing a single component---directly supports the operational requirements of the platform.


\Cref{sec:irie} presents the implementation of a structural health monitoring framework for the \BRACE{} project.
Its computational backend, the Infrastructure Resilience Engine (IRiE), organizes health monitoring around five abstractions: \emph{Assets} representing physical structures, \emph{Events} representing earthquakes or other triggering conditions, \emph{Predictors} implementing physics-based or data-driven analysis methods, \emph{Metrics} quantifying structural condition, and \emph{Evaluations} assembling metrics into actionable reports.
This architecture permits physics-based simulation and system identification to coexist within a unified framework, with results presented through a common interface. 
The developments of Part~I contribute the simulation engine for Type~I predictors on this platform.
Following a pilot deployment covering 22 bridges, Caltrans has funded continued development to extend coverage across the California highway network.

\Cref{ch:predictor_t1}

\Cref{ch:xara} develops a software architecture for finite element simulations that addresses the demands of \cref{sec:irie}. The developments are implemented in the open source project \Xara{}, which drives the physics-based simulations on the \BRACE{} platform. 

\Cref{ch:autodiff}
